\frontsection{Современные модели для таблиц: дополнительное чтение}
\label{sec:modern-tabular}
\subsection*{От изученной сети к новым подходам}
Во второй работе полносвязная сеть обучалась на признаках домов. Её также называют MLP (multilayer perceptron, многослойный перцептрон). Исследователи продолжают развивать такие модели: меняют способ объединения прогнозов и используют предварительное обучение. Ниже рассмотрены две идеи. Их реализация не требуется для выполнения основных работ.

\subsection*{TabM: несколько прогнозов внутри одной модели}
\textbf{Ансамбль} объединяет ответы нескольких моделей. Представим три оценки цены одного дома: 180, 190 и 200 условных единиц. Их среднее равно 190. Если ошибки участников различаются, усреднение может уменьшить ошибку; одинаковые систематические промахи таким способом не устраняются.

Обычный ансамбль отдельных нейросетей требует обучить и хранить каждую из них. В \textbf{TabM}, представленном на ICLR 2025, несколько вариантов MLP совместно используют большую часть параметров и обучаются вместе. Модель формирует несколько прогнозов для объекта и объединяет их. Это позволяет экономнее реализовать идею ансамбля~\cite{tabm}.

Для первого знакомства достаточно различать «одна сеть --- один прогноз» и «несколько согласованно обучаемых вариантов --- объединённый прогноз». Приведённое выше среднее иллюстрирует объединение ответов, но не воспроизводит устройство TabM.

\subsection*{TabPFN-3: использовать опыт предварительного обучения}
\textbf{Предобученной} называют модель, параметры которой настроены заранее. В \textbf{TabPFN-3} предварительное обучение выполняется на искусственно созданных таблицах. При решении новой задачи модель получает примеры с известными ответами и признаки объектов, для которых ответ нужно предсказать. Обучающие примеры выступают контекстом для прогноза; это называют \textit{in-context learning} --- обучением на примерах в контексте~\cite{tabpfn3}.

Здесь важно различать два этапа: затратное предварительное обучение и применение готовой модели к новой таблице. Во втором случае прогноз не обязательно требует обычной настройки весов сети на каждом новом наборе. Модель при этом всё равно использует известные ответы обучающих примеров.

Источник по TabPFN-3 --- \textbf{технический отчёт мая 2026 года}. Заявления о качестве относятся к экспериментам его авторов. Они не устанавливают качество на файлах нашего курса и не означают, что метод будет лучшим для любой таблицы.

\clearpage
\subsection*{Что сравнивать, кроме названия модели}
Новый метод оценивают по тем же вопросам, что и изученную MLP: какие данные он получает, как разделены примеры, какая ошибка измеряется, сколько времени и памяти требуется. Предобучение --- дополнительный ресурс, который нужно указать. Улучшение метрики на чужом наборе не заменяет проверку на своей задаче.

\begin{table}[H]\centering\small
\begin{tabularx}{\textwidth}{@{}p{0.22\textwidth}Y@{}}
\toprule
Подход & Основная идея \\
\midrule
MLP из работы №\,2 & Настроить одну полносвязную сеть на обучающих домах. \\
TabM & Совместно обучать несколько вариантов предсказания с общими параметрами. \\
TabPFN-3 & Применить заранее обученную модель, используя примеры новой таблицы как контекст. \\
\bottomrule
\end{tabularx}
\end{table}

\subsection*{Необязательное задание: помогает ли усреднение?}
Этот опыт можно выполнить средствами работы №\,2, не устанавливая TabM или TabPFN-3. Он знакомит с идеей ансамбля отдельных сетей; это не реализация этих методов.
\begin{enumerate}
\item Сохраните одно разбиение домов и одну предобработку из работы №\,2. Выберите исходную конфигурацию B0.
\item Обучите три новые копии этой сети с seed 40, 41 и 42. Seed обучения влияет на начальные веса и случайный порядок объектов. Разбиение данных и остальные настройки оставьте теми же.
\item Для каждой сети сохраните прогнозы валидации в исходных единицах цены. В учебном модуле это поле \file{validation\_predictions} результата \file{run\_experiment}.
\item Для каждого дома усредните три прогноза:
\[
\widehat y_i^{\,avg}=
\frac{\widehat y_i^{(1)}+\widehat y_i^{(2)}+\widehat y_i^{(3)}}{3}.
\]
Усредняются ответы для одного и того же дома, а не значения MAE трёх запусков.
\item Составьте четыре строки результата: три отдельные сети и усреднение. Для каждой укажите MAE и RMSE на одной валидации. Сравните также затраты на три обучения и на одно.
\end{enumerate}

Покажите один дом, для которого усреднение помогло, и один, для которого не помогло, если такие примеры есть. Объясните наблюдение. Укажите, что преимущество на одном разбиении не доказывает общего превосходства, а разброс трёх прогнозов сам по себе не является гарантированным интервалом цены.

\textbf{Форма сдачи:} небольшая таблица, два разобранных примера и вывод на полстраницы. Задание дополнительное и выполняется в среде работы №\,2. Для него не требуются библиотеки TabM/TabPFN-3, доступ к их API или их предобученные веса.
